\section[D. 左灯右行]{\myboxiv{      \BF{\TT{D}}} 左灯右行}

\textit{运行时间限制}   1 second

\textit{内存使用限制} 1 Gigabyte

\textit{读入文件} \qquad 标准输入

\textit{输出文件} \qquad 标准输出


\begin{quote}
	\textit{我爱你有种【数据删除】的冲突，}
\end{quote}

\BF{本题和试题 \myboxv{\BF{\TT{E}}} 右灯左行共用相同的背景，两道题之间有较大的关系，建议在尝试其中一题前首先阅读另一题的文本。}

\BF{这是一道交互题。}

\subsection{题目描述}

小 T 面对 T 城的地铁图犯了难 ...

具体而言，这座城市在 E 河两岸，北岸有 $n$ 个地铁站点 $N_1$, $N_2$, $\ldots$, $N_n$, 南岸有 $n$ 个地铁站点 $S_1$, $S_2$, $\ldots$, $S_n$.

T 城总共有 $n - 1$ 条地铁线路，每条线路都是环线，其中第 $i$ 条环线连接的地铁站是 $N_i - N_{i + 1} - S_{i + 1} - S_i - N_i$, 但这些地铁环线都是单向运行的，可能是 $N_i \leftarrow N_{i + 1} \leftarrow S_{i + 1} \leftarrow S_i \leftarrow N_i$ 也可能是 $N_i \rightarrow N_{i + 1} \rightarrow S_{i + 1} \rightarrow S_i \rightarrow N_i$.


每条环线都有一个价目表，分别表示这 4 个相邻站 ($N_i - N_{i + 1}$, $N_{i + 1} - S_i$, $S_i - S_{i + 1}$, $N_i - S_i$) 之间坐一站地铁移动的价格，注意坐一站地铁的方向由地铁的运行方向而定，可能是 $N_i \rightarrow N_{i + 1}$ 也可能是 $N_{i} \leftarrow N_{i + 1}$.

遗憾的是，时空黑洞吞噬了价目表上这些环线的运行方向和价格，小 T 准备求助大 K 的帮助来得知这些地铁的运行方向，但作为回报，小 T 需要帮助大 K 解决一系列问题。

在第一阶段，大 K 作为另一个时空的人，他知道地铁的运行方向，小 T 需要首先指定地铁图上的每条环线的每个相邻站之间的价目表，随后根据指定的价目表询问大 K 得到地铁的运行方向。

小 T 可以询问他 $q$ 个问题：

\begin{itemize}
	\item 城市 $x$(可以是 $N_i$ 也可以是 $S_i$, 下略) 到城市 $y$ 坐地铁，在小 T 指定的地铁价目表下，至少需要支付多少地铁费用？
\end{itemize}

在这个问题下，$Q = 1\,500$.

在询问完这个问题后，小 T 需要和大 K 确认每条地铁的运行方向，随后的第二阶段，大 K 将会反过来询问小 T $q$ 个问题：

\begin{itemize}
	\item 城市 $x$(可以是 $N_i$ 也可以是 $S_i$, 下略) 到城市 $y$ 坐地铁，在小 T 指定的地铁价目表下，至少需要支付多少地铁费用？
\end{itemize}

请帮助小 T 解决这两个阶段的问题。

请注意，地铁线路的方向是预先确定的，他们不会在交互的过程中更改，换言之交互器\BF{不是}自适应的。

\subsection{交互协议}

本题每个测试点只有一组测试数据。

第一行输入一行两个整数 $n$ 和 $q$ ($1 \le n, q \le 5 \times 10^4$), 表示城市数量和地铁数量的参数。

随后进入交互第一阶段。

在第一阶段，你首先需要输出 $n - 1$ 行，其中第 $i$ 行输出 $4$ 个整数 $u_i$, $d_i$, $l_i$, $r_i$, 分别表示第 $i$ 条环线上 $N_i - N_{i + 1}$, $S_i - S_{i + 1}$, $N_i - S_i$ 和 $N_{i + 1} - S_{i + 1}$ 这四个相邻站之间的移动价格，你需要保证这四个整数在 $[0, 10^9]$ 之间。

随后你将会直接开始交互，每一次交互你可以以下面的格式向交互器询问问题：

\begin{itemize}
   \item 输出一行，首先输出一个字符 \TT{?} 表示你询问一次，随后输出四个元素 $S_x$, $id_x$, $S_y$, $id_y$, 其中你需要保证 $S_x, S_y \in \{\mathtt{N}, \mathtt{S}\}$, $1 \le id_x, id_y \le n$, 表示小 T 需要询问大 K 位于 $S_x$ 岸的第 $id_x$ 个城市 $(S_x)_{id_x}$ 到位于 $S_y$ 岸的第 $id_y$ 个城市 $(S_y)_{id_y}$ 之间至少需要支付多少地铁费用。
   \item 若你的输入合法并且没有超出询问限制，交互器会给出一行一个整数 $d$ 供输入表示这两个站点之间最少需要的地铁费用。
   \item 若你询问的问题超出了对应的询问次数限制，交互器会输出一个整数 $-1$ 并结束你的交互进程，此时你会获得\textit{答案错误}的评测结果。
\end{itemize}

若你已经得到地铁方向的话，请按照以下的格式输出答案：

\begin{itemize}
   \item 输出一行，首先输出一个字符 \TT{!}.
   \item 随后输出一个长度为 $n - 1$ 的字符串 $S$, 其中第 $i$ 个字符 $S_i \in \{\mathtt{I}, \mathtt{O}\}$, 若 $S_i = \mathtt{I}$ 则表示第 $i$ 条环线的方向为 $N_i \leftarrow N_{i + 1} \leftarrow S_{i + 1} \leftarrow S_i \leftarrow N_i$, 若 $S_i = \mathtt{O}$ 则表示第 $i$ 条环线的方向为 $N_i \rightarrow N_{i + 1} \rightarrow S_{i + 1} \rightarrow S_i \rightarrow N_i$.
\end{itemize}

随后将会进入第二阶段，交互器会询问你 $Q$ 个问题，格式如下：

\begin{itemize}
   \item 输入一行四个元素 $S_x$, $id_x$, $S_y$, $id_y$, ($S_x, S_y \in \{\mathtt{N}, \mathtt{S}\}$, $1 \le id_x, id_y \le n$), 表示大 K 询问小 T 位于 $S_x$ 岸的第 $id_x$ 个城市 $(S_x)_{id_x}$ 到位于 $S_y$ 岸的第 $id_y$ 个城市 $(S_y)_{id_y}$ 之间至少需要支付多少地铁费用。
   \item 你需要输出一行一个整数 $d$ 表示这两个点之间最少需要的地铁费用。
\end{itemize}

在完成一组交互之后你应该立刻开始下一组交互的询问，若已是最后一组数据你应当安全退出。

每次输出都需要输出换行并且刷新缓冲区，否则你可能会得到意料之外的结果。

为了刷新缓冲区，你可以：

\begin{itemize}
   \item 对于 C 或者 C++, 使用 \TT{fflush(stdout)} 或 \TT{cout.flush()}。
   \item 对于 Java 或 Kotlin, 使用 \TT{System.out.flush()}。
   \item 对于 Python, 使用 \TT{stdout.flush()}。
\end{itemize}

\subsection{样例 1}

\begin{minipage}{\linewidth}
	\begin{minipage}[t]{0.49\linewidth}
		\begin{lstlisting}[title=\texttt{standard input}]
TBD
\end{lstlisting}
	\end{minipage}
	\begin{minipage}[t]{0.49\linewidth}
		\begin{lstlisting}[title=\texttt{standard output}]
TBD
\end{lstlisting}
	\end{minipage}
\end{minipage}


\subsection{样例 1 解释}

请注意样例中的交互过程仅供参考，实际交互过程并不唯一，且这一示例的交互过程不一定是可行的或是最优的，本题的样例不会在附加文件中出现。





